\chapter{Metodologia}
\label{chap:methodology}

Este capítulo descreve a abordagem metodológica adotada
no desenvolvimento da Medicano, abrangendo o tipo de
pesquisa, o método de desenvolvimento, as ferramentas
utilizadas, o ambiente de desenvolvimento e a estratégia
de testes.

\section{Tipo de Pesquisa}
\label{sec:research-type}

O presente trabalho se caracteriza como uma pesquisa
aplicada, uma vez que tem como objetivo produzir um
artefato de software funcional para solucionar um
problema concreto identificado no setor de saúde
\cite{gil2002}. Do ponto de vista dos procedimentos
técnicos, trata-se de uma pesquisa de desenvolvimento,
na qual os conhecimentos adquiridos ao longo do curso
de Engenharia de Software são aplicados na construção
de um sistema real com propósito definido.

A abordagem é predominantemente prática e iterativa,
alinhada com os princípios do desenvolvimento ágil,
nos quais o aprendizado ocorre de forma contínua ao
longo dos ciclos de desenvolvimento e não apenas ao
final do projeto \cite{schwaber2020}.

\section{Método de Desenvolvimento}
\label{sec:dev-method}

O desenvolvimento da Medicano combina dois métodos
complementares: o \textit{Spec Driven Development}
(SDD) como abordagem de especificação e o Scrum como
framework de execução iterativa. A execução das etapas
de codificação a partir das specs aprovadas é apoiada
por um pipeline de agentes de IA orquestrados pelo
CrewAI, descrito na Seção~\ref{subsec:ai-assisted-dev}.

\subsection{Spec Driven Development}
\label{subsec:sdd}

O SDD estabelece que toda funcionalidade deve ser
completamente especificada antes de sua implementação.
A especificação define contratos precisos — modelos
de dados, regras de negócio, fluxos e interfaces da
API — que orientam o desenvolvimento e servem de
critério de aceitação ao final de cada sprint.

Na Medicano, as especificações são elaboradas e
versionadas em um diretório \texttt{specs/} no próprio
repositório de código, seguindo um \textit{template}
interno que padroniza as seções de cada \textit{spec}:
contexto, requisitos funcionais e regras de negócio
aplicáveis, entidades envolvidas, decisões técnicas
não-negociáveis e critérios de aceite. A cada
funcionalidade, uma nova \textit{spec} é aberta em uma
\textit{branch} dedicada, submetida à revisão por pares
em Pull Request e, após aprovação, serve de contrato
para a etapa de implementação.
O Capítulo~\ref{chap:specification} é o produto
direto dessa abordagem: toda funcionalidade
implementada passou por um ciclo de
especificação~$\rightarrow$~revisão~$\rightarrow$~implementação.

O SDD traz dois benefícios centrais para o projeto:
reduz retrabalho ao detectar inconsistências de
requisito antes do código ser escrito, e facilita
a divisão de trabalho entre os membros da equipe,
que podem implementar em paralelo a partir de um
contrato comum.

\subsection{Scrum}
\label{subsec:scrum}

O Scrum organiza a execução do trabalho em sprints
quinzenais \cite{schwaber2020}. Cada sprint inicia
com a revisão das specs aprovadas e termina com uma
entrega funcional validada. Essa cadência permite
inspeção regular do progresso e ajuste de prioridades
sem comprometer o planejamento geral.

O backlog do produto é organizado por ordem de
prioridade de negócio, sendo os fluxos essenciais do
MVP — autenticação, triagem por chatbot, busca e
agendamento — priorizados nas primeiras sprints. A
Tabela~\ref{tab:sprints} apresenta a distribuição do
escopo ao longo das sprints planejadas.

\begin{table}[H]
  \centering
  \caption{Planejamento de sprints}
  \label{tab:sprints}
  \begin{tabular}{cp{6cm}p{3.5cm}p{2cm}}
    \toprule
    \textbf{Sprint} &
    \textbf{Escopo} &
    \textbf{Entregável} &
    \textbf{Duração} \\
    \midrule
    1 & Setup do monorepo, configuração AWS,
        modelos de dados, autenticação
      & Auth funcionando
      & 2 semanas \\
    \midrule
    2 & Módulo de clínicas e profissionais,
        busca com filtros, perfil público
      & Busca de profissionais
      & 2 semanas \\
    \midrule
    3 & Chatbot de triagem com Vercel AI SDK,
        integração com busca da plataforma
      & Triagem funcionando
      & 2 semanas \\
    \midrule
    4 & Sistema de agendamentos, regras de
        negócio, validações e cancelamento
      & Agendamento \textit{end-to-end}
      & 2 semanas \\
    \midrule
    5 & Dashboard da clínica, gestão de agenda,
        confirmação e cancelamento
      & Painel da clínica
      & 2 semanas \\
    \midrule
    6 & Módulo de assinaturas, planos de clínica
        e profissional autônomo, integração de
        pagamento
      & Assinaturas funcionando
      & 2 semanas \\
    \midrule
    7 & Testes de aceitação, ajustes de UX,
        otimizações de desempenho, deploy
        final na AWS
      & Sistema completo em produção
      & 2 semanas \\
    \bottomrule
  \end{tabular}
  \fonte{Elaborado pelos autores (2026)}
\end{table}

\subsection{Desenvolvimento Assistido por IA}
\label{subsec:ai-assisted-dev}

O desenvolvimento da Medicano é apoiado por um pipeline
de agentes de inteligência artificial orquestrados pelo
framework CrewAI. Essa escolha substitui o uso de
assistentes de IA genéricos --- como ferramentas de
preenchimento automático ou conversação ad hoc com um
único modelo --- por um fluxo estruturado em que cada
etapa da construção de uma funcionalidade é conduzida
por um agente especializado, com papel, \textit{prompt}
e modelo de linguagem próprios.

O pipeline é disparado uma vez por funcionalidade, a
partir da \textit{spec} previamente aprovada no
diretório \texttt{specs/} do repositório. A
\textit{spec} é fornecida como entrada do
primeiro agente, e o resultado de cada etapa é
encadeado à próxima em um processo sequencial que
reproduz o ciclo de especificação~$\rightarrow$~implementação~$\rightarrow$~documentação~$\rightarrow$~revisão
adotado pelo projeto. A
Tabela~\ref{tab:ai-agents} descreve os quatro agentes
que compõem o pipeline, sua responsabilidade no ciclo
e o modelo de linguagem utilizado por cada um.

\begin{table}[H]
  \centering
  \caption{Agentes de IA do pipeline de desenvolvimento}
  \label{tab:ai-agents}
  \begin{tabular}{p{2.6cm}p{5.2cm}p{5.6cm}}
    \toprule
    \textbf{Agente} &
    \textbf{Responsabilidade} &
    \textbf{Modelo de linguagem} \\
    \midrule
    Arquiteto
      & Converter a \textit{spec} aprovada em plano
        técnico contendo arquivos a serem criados ou
        modificados, DTOs, \textit{endpoints} e regras
        de negócio a validar
      & OpenAI o4-mini (via OpenRouter) \\
    \midrule
    Desenvolvedor
      & Implementar o código TypeScript descrito no
        plano seguindo as convenções do monorepo
        (NestJS, Mongoose e React)
      & Anthropic Claude Sonnet 4.5
        (via OpenRouter) \\
    \midrule
    Documentador
      & Anotar o código com JSDoc e decorators do
        Swagger e produzir \textit{changelog}
        rastreável aos RF, RN e UC correspondentes
      & Google Gemini 2.5 Flash
        (via OpenRouter) \\
    \midrule
    Revisor
      & Verificar segurança, aderência às regras de
        negócio e conformidade com LGPD e CFM
        antes da integração à branch
        \texttt{main}
      & OpenAI o4-mini (via OpenRouter) \\
    \bottomrule
  \end{tabular}
  \fonte{Elaborado pelos autores (2026)}
\end{table}

A integração entre provedores distintos é feita pelo
OpenRouter, um \textit{gateway} que expõe uma interface
unificada para modelos de diferentes fornecedores. Essa
abstração permite alternar o modelo de cada agente sem
alterar o código do \textit{pipeline} e reduz o impacto
de indisponibilidades pontuais de um provedor
específico. Adicionalmente, o Claude é utilizado
diretamente, fora do \textit{pipeline} do CrewAI, para
apoio em decisões de arquitetura de mais alto nível e
na elaboração desta documentação.

A decisão final sobre qualquer alteração no código-fonte
continua sendo humana. O resultado do \textit{pipeline}
é submetido à mesma revisão por par em Pull Request
descrita na Seção~\ref{subsec:version-control} antes de
ser integrado à branch principal, garantindo que nenhum
código gerado por IA chegue à \texttt{main} sem
avaliação de ao menos um integrante da equipe.

\section{Ferramentas Utilizadas}
\label{sec:tools}

As ferramentas utilizadas no desenvolvimento da Medicano
foram selecionadas com base na compatibilidade com a
stack tecnológica do projeto e na adoção pela comunidade
de desenvolvimento de software. A
Tabela~\ref{tab:tools} apresenta as principais
ferramentas e suas finalidades.

\begin{table}[H]
  \centering
  \caption{Ferramentas de desenvolvimento}
  \label{tab:tools}
  \begin{tabular}{p{3.8cm}p{3cm}p{7cm}}
    \toprule
    \textbf{Ferramenta} &
    \textbf{Categoria} &
    \textbf{Finalidade} \\
    \midrule
    Visual Studio Code  & Editor de código
      & Desenvolvimento do frontend e backend \\
    Git                 & Controle de versão
      & Versionamento do código-fonte \\
    GitHub              & Repositório remoto
      & Hospedagem, colaboração e gestão de specs \\
    CrewAI              & Framework de agentes de IA
      & Orquestração do \textit{pipeline} de
        desenvolvimento assistido por IA descrito
        na Seção~\ref{subsec:ai-assisted-dev} \\
    OpenRouter          & \textit{Gateway} de modelos de IA
      & Interface unificada para os modelos de
        linguagem utilizados pelos agentes do
        CrewAI \\
    Claude              & Assistente de IA
      & Apoio em decisões de arquitetura, revisão
        de código e elaboração da documentação \\
    Jira                & Gestão de projeto
      & Controle de backlog, sprints e issues \\
    Confluence          & Documentação interna
      & Base de conhecimento e onboarding \\
    Jira Service Management & Suporte ao cliente
      & Gerenciamento de chamados e atendimento \\
    Insomnia            & Cliente HTTP
      & Testes manuais da API \\
    MongoDB Compass     & Interface de banco
      & Visualização e consulta do MongoDB \\
    Overleaf            & Editor \LaTeX
      & Elaboração desta documentação \\
    Lucidchart          & Diagramação
      & Criação dos diagramas de arquitetura
        e fluxo do sistema \\
    Figma               & Design de interface
      & Prototipação das telas do sistema \\
    AWS Console         & Painel de controle
      & Gerenciamento da infraestrutura AWS \\
    Docker              & Contêinerização
      & Serviços locais de banco de dados \\
    Docker Compose      & Orquestração local
      & Subida de MongoDB e Redis em desenvolvimento \\
    \bottomrule
  \end{tabular}
  \fonte{Elaborado pelos autores (2026)}
\end{table}

\subsection{Controle de Versão com Git e GitHub}
\label{subsec:version-control}

O controle de versão do código-fonte é realizado com
Git, utilizando o GitHub como repositório remoto
central \cite{chacon2014}. O Git permite que a equipe
trabalhe de forma paralela em diferentes funcionalidades
sem conflitos, mantendo um histórico completo de todas
as alterações realizadas no projeto.

O repositório é organizado como um monorepo, contendo
tanto o frontend quanto o backend em uma única estrutura
versionada. Essa abordagem facilita o desenvolvimento
coordenado entre as duas camadas, uma vez que alterações
que afetam simultaneamente o frontend e o backend podem
ser versionadas em um único commit, mantendo a
consistência do projeto.

O fluxo de trabalho adotado é baseado no
\textit{GitHub Flow}, uma estratégia que organiza o
desenvolvimento em torno de uma branch principal estável
e branches de funcionalidade \cite{chacon2014}. O
processo funciona da seguinte forma:

\begin{itemize}
  \item A branch \texttt{main} representa sempre o
        estado estável e deployável do projeto
  \item Cada nova funcionalidade ou correção é
        desenvolvida em uma branch separada, nomeada
        de forma descritiva — por exemplo,
        \texttt{feature/chat-triagem} ou
        \texttt{fix/slot-conflict}
  \item Ao finalizar o desenvolvimento, o membro da
        equipe abre um Pull Request no GitHub para
        revisão do código por outro integrante
  \item Após aprovação, a branch é integrada à
        \texttt{main} e o deploy pode ser realizado
\end{itemize}

Essa prática garante que nenhum código não revisado
chegue à branch principal, reduzindo a incidência de
bugs em produção e promovendo a transferência de
conhecimento entre os membros da equipe por meio das
revisões de código \cite{chacon2014}.

\subsection{Gestão do Projeto com Atlassian}
\label{subsec:atlassian}

O gerenciamento do projeto é realizado por meio do
ecossistema Atlassian, composto pelo Jira, Confluence
e Jira Service Management — ferramentas amplamente
adotadas pela indústria de software para gestão de
projetos ágeis e suporte a produtos digitais
\cite{atlassian2024}.

O \textbf{Jira} é utilizado como ferramenta central
de gestão do backlog e das sprints. Cada funcionalidade
do sistema é registrada como uma \textit{user story},
estimada em story points e priorizada conforme o valor
de negócio. As sprints quinzenais são criadas e
acompanhadas no Jira, permitindo que a equipe visualize
o progresso por meio de boards Kanban e gráficos de
burndown.

O \textbf{Confluence} serve como base de conhecimento
operacional da equipe, centralizando toda a informação
que um novo colaborador precisa para começar a
contribuir com o projeto. Isso inclui guias de
configuração do ambiente de desenvolvimento, decisões
de arquitetura registradas como ADRs (Architecture
Decision Records), padrões de código adotados pelo
time, fluxos operacionais e atas das reuniões de
sprint. A integração nativa com o Jira permite
referenciar issues e sprints diretamente nas páginas
do Confluence, mantendo a rastreabilidade entre
decisões documentadas e tarefas implementadas.

O \textbf{Jira Service Management} é utilizado para
o gerenciamento do suporte ao cliente da plataforma.
Ele permite que usuários da Medicano registrem chamados
diretamente pela interface da plataforma, os quais são
automaticamente criados como tickets e roteados para
a equipe responsável. A integração com o Jira de
desenvolvimento permite que bugs reportados pelos
usuários sejam escalados diretamente para o backlog
do produto, garantindo rastreabilidade entre o chamado
do cliente e a correção implementada.

\section{Ambiente de Desenvolvimento}
\label{sec:dev-environment}

O ambiente de desenvolvimento local é configurado de
forma padronizada para todos os membros da equipe,
garantindo consistência entre as máquinas e evitando
divergências de comportamento entre ambientes. A
padronização é realizada por meio de arquivos de
configuração versionados no repositório, incluindo
\texttt{.nvmrc} para fixar a versão do Node.js e
\texttt{.eslintrc} e \texttt{.prettierrc} para
padronização do estilo de código.

As variáveis de ambiente sensíveis são gerenciadas
exclusivamente por meio do AWS Secrets Manager em
todos os ambientes, sem o uso de arquivos \texttt{.env}
locais \cite{awssm}. Essa decisão elimina o risco de
credenciais sensíveis serem acidentalmente versionadas
no repositório e garante consistência no mecanismo de
acesso a segredos entre desenvolvimento e produção.

A distinção entre os três ambientes é realizada por
meio de prefixos de secrets diferentes no Secrets
Manager:

\begin{itemize}
  \item \textbf{Desenvolvimento}: secrets prefixados
        com \texttt{medicano/development/}, conectando
        ao database \texttt{medicano\_development} do
        cluster Atlas
  \item \textbf{\textit{Staging}}: secrets prefixados
        com \texttt{medicano/staging/}, conectando ao
        database \texttt{medicano\_staging} — replica
        a configuração de produção para validação
        antes da publicação
  \item \textbf{Produção}: secrets prefixados com
        \texttt{medicano/production/}, conectando ao
        database \texttt{medicano\_production}
\end{itemize}

A aplicação identifica o ambiente ativo por meio da
variável \texttt{NODE\_ENV}, que é a única variável
definida localmente e não contém informações sensíveis.
Com base nesse valor, o serviço de configuração do
NestJS determina qual prefixo de secret utilizar na
inicialização da aplicação.

O banco de dados MongoDB e o Redis são executados
localmente em contêineres Docker, garantindo que
todos os membros da equipe utilizem as mesmas versões
e configurações \cite{docker2024}. O ambiente é
inicializado com o comando
\mbox{\texttt{docker-compose up -d}} na raiz do
repositório, subindo os serviços em contêineres
isolados sem interferir em outros projetos da máquina.

Para os testes de integração, o MongoDB é
substituído por uma instância em memória provida
pelo \texttt{mongodb-memory-server}, que inicializa
um servidor MongoDB temporário durante a execução
dos testes e o encerra ao final, garantindo
isolamento total e reprodutibilidade dos resultados
\cite{jestjs2024}. O Redis é substituído pelo
\texttt{ioredis-mock}, uma implementação em memória
compatível com a interface do cliente Redis real.

A Tabela~\ref{tab:environment} descreve as principais
tecnologias e versões utilizadas no ambiente de
desenvolvimento.

\begin{table}[H]
  \centering
  \caption{Ambiente de desenvolvimento}
  \label{tab:environment}
  \begin{tabular}{lll}
    \toprule
    \textbf{Tecnologia} &
    \textbf{Versão} &
    \textbf{Finalidade} \\
    \midrule
    Node.js          & 20 LTS   & Execução do backend \\
    npm              & 10+      & Gerenciamento de pacotes \\
    TypeScript       & 5+       & Linguagem principal \\
    NestJS           & 10+      & Framework backend \\
    React            & 18+      & Framework frontend \\
    Vite             & 5+       & Build tool do frontend \\
    Docker           & 24+      & Contêinerização dos serviços \\
    Docker Compose   & 2+       & Orquestração do ambiente local \\
    MongoDB          & 7+       & Banco de dados (via Docker) \\
    Redis            & 7+       & Gerenciamento de tokens (via Docker) \\
    \bottomrule
  \end{tabular}
  \fonte{Elaborado pelos autores (2026)}
\end{table}

\section{Estratégia de Testes}
\label{sec:testing-strategy}

A estratégia de testes da Medicano é organizada em
três níveis complementares, cada um com escopo e
objetivo distintos \cite{sommerville2011}.

\textbf{Testes unitários} são aplicados à camada de
services do backend, onde está concentrada a lógica
de negócio da aplicação. Cada service é testado de
forma isolada, com as dependências externas — como
o banco de dados e as APIs de terceiros —
substituídas por mocks. O framework utilizado é o
Jest, já integrado ao NestJS em sua configuração
padrão \cite{jestjs2024}.

\textbf{Testes de integração} verificam o
comportamento dos endpoints da API de ponta a ponta,
desde o recebimento da requisição até a resposta
retornada ao cliente. Esses testes utilizam um banco
de dados MongoDB em memória para evitar dependência
de infraestrutura externa durante a execução.

\textbf{Testes manuais} são realizados com o
Insomnia para validar os fluxos completos da API
durante o desenvolvimento, e com navegadores para
validação da interface do usuário antes de cada
entrega de sprint.

A cobertura de testes automatizados é focada nos
fluxos críticos do sistema: autenticação, regras de
agendamento e lógica de triagem. Fluxos secundários
são cobertos pelos testes manuais realizados ao
final de cada sprint, conforme o processo definido
pelo Scrum \cite{schwaber2020}.